% Part: lambda-calculus
% Chapter: syntax
% Section: alpha

\documentclass[../../../include/open-logic-section]{subfiles}

\begin{document}

\olfileid[id]{lam}{syn}{alp}
\olsection{Konversi-$\alpha$}

Apakah hubungan antara $\lambd[x][x]$ dan $\lambd[y][y]$? Keduanya
menyatakan fungsi identitas. Tentu saja, secara sintaktis keduanya merupakan
suku yang berbeda. Perbedaannya hanya terletak pada nama variabel terikat, dan
salah satunya merupakan hasil ``penggantian nama'' variabel terikat dalam suku
yang lain. Hal ini disebut \emph{konversi-$\alpha$}.

\begin{defn}[Pengubahan variabel terikat, $\aconvone$]
  Jika suatu suku $M$ memuat kemunculan $\lambd[x][N]$, $x \neq y$,
  $y \notin \FV{N}$, dan $\Subst{N}{y}{x}$ terdefinisi, maka penggantian
  kemunculan ini dengan
  \begin{equation*}
    \lambd[y][\Subst{N}{y}{x}]
  \end{equation*}
  yang menghasilkan $M'$ disebut \emph{pengubahan variabel terikat}, dan
  ditulis $M \redone[\alpha] M'$.
\end{defn}

\begin{defn}[Kompatibilitas relasi]
  Suatu relasi $R$ pada suku-suku dikatakan \emph{kompatibel}
  jika memenuhi syarat-syarat berikut:
  \begin{enumerate}
  \item Jika $R N N'$, maka $R \lambd[x][N] \lambd[x][N']$.
  \item Jika $R P P'$, maka $R (PQ) (P'Q)$.
  \item Jika $R Q Q'$, maka $R (PQ) (PQ')$.
  \end{enumerate}
\end{defn}

Jadi, mari kita rumuskan ulang definisi tersebut:
\begin{defn}[Pengubahan variabel terikat, $\aconvone$]
  \emph{Pengubahan variabel terikat} ($\redone[\alpha]$) adalah
  relasi kompatibel terkecil pada suku-suku yang memenuhi syarat berikut:
  \begin{align*}
    &\lambd[x][N] \redone[\alpha] \lambd[y][\Subst{N}{y}{x}] && \text{jika
      $x \neq y$, $y \notin \FV{N}$} \\
    & &&\text{dan $\Subst{N}{y}{x}$ terdefinisi}
  \end{align*}
\end{defn}

``Terkecil'' di sini berarti bahwa relasi tersebut hanya memuat pasangan yang
disyaratkan oleh kompatibilitas dan syarat tambahan di atas, dan tidak memuat
pasangan lain. Jadi, relasi ini juga dapat didefinisikan sebagai berikut:

\begin{defn}[Pengubahan variabel terikat, $\aconvone$] \ollabel{defn:aconvone}
  \emph{Pengubahan variabel terikat} ($\aconvone$) didefinisikan secara
  induktif sebagai berikut:
  \begin{enumerate}
  \item Jika $N \aconvone N'$, maka $\lambd[x][N] \aconvone
    \lambd[x][N']$. \ollabel{defn:aconvone1}
  \item Jika $P \aconvone P'$, maka $(PQ) \aconvone (P'Q)$.
    \ollabel{defn:aconvone2}
  \item Jika $Q \aconvone Q'$, maka $(PQ) \aconvone (PQ')$.
    \ollabel{defn:aconvone3}
  \item Jika $x \neq y$, $y \notin \FV{N}$, dan $\Subst{N}{y}{x}$
    terdefinisi, maka
    $\lambd[x][N] \redone[\alpha] \lambd[y][\Subst{N}{y}{x}]$.
    \ollabel{defn:aconvone4}
  \end{enumerate}
\end{defn}

Definisi-definisi tersebut ekuivalen, tetapi pembuktiannya kita serahkan
sebagai latihan. Mulai sekarang, kita akan menggunakan definisi induktif.

\begin{defn}[Konversi-$\alpha$, $\aconv$]
  \emph{Konversi-$\alpha$} ($\aconv$) adalah relasi refleksif dan transitif
  terkecil pada suku-suku yang memuat $\aconvone$.
\end{defn}

Seperti di atas, ``terkecil'' berarti bahwa relasi tersebut hanya memuat
pasangan yang disyaratkan oleh refleksivitas, transitivitas, dan $\aconvone$,
sehingga kita memperoleh definisi ekuivalen berikut:

\begin{defn}[Konversi-$\alpha$, $\aconv$] \ollabel{defn:aconv}
  \emph{Konversi-$\alpha$} ($\aconv$) didefinisikan secara induktif sebagai
  berikut:
  \begin{enumerate}
  \item Jika $P \aconv Q$ dan $Q \aconv R$, maka $P \aconv R$.
    \ollabel{defn:aconv1}
  \item Jika $P \aconvone Q$, maka $P \aconv Q$. \ollabel{defn:aconv2}
  \item $P \aconv P$. \ollabel{defn:aconv3}
  \end{enumerate}
\end{defn}

\begin{ex}
  $\lambd[x][f x]$ dapat dikonversi-$\alpha$ menjadi $\lambd[y][f y]$, dan
  sebaliknya. Secara informal, keduanya merupakan fungsi yang menerima sebuah
  argumen dan mengembalikan hasil penerapan $f$ pada argumen tersebut, dengan
  $f$ merujuk kepada variabel lingkungan.

  $\lambd[x][f x]$ tidak dapat dikonversi-$\alpha$ menjadi $\lambd[x][g x]$.
  Secara informal, kedua suku itu masing-masing merujuk kepada variabel
  lingkungan $f$ dan~$g$, sehingga merupakan fungsi yang berbeda: keduanya
  berperilaku berbeda dalam lingkungan tempat $f$ dan $g$ berbeda.
\end{ex}

\begin{prob}
  Apakah pasangan-pasangan suku berikut dapat dikonversi-$\alpha$?
  \begin{enumerate}
  \item $\lambd[x][\lambd[y][x]]$ dan $\lambd[y][\lambd[x][y]]$
  \item $\lambd[x][\lambd[y][x]]$ dan $\lambd[c][\lambd[b][a]]$
  \item $\lambd[x][\lambd[y][x]]$ dan $\lambd[c][\lambd[b][a]]$
  \end{enumerate}
\end{prob}

\begin{lem}\ollabel{lem:fv-one}
  Jika $P \aconvone Q$, maka $\FV{P} = \FV{Q}$.
\end{lem}

\begin{proof}
  Dengan induksi pada !!{derivation} $P \aconvone Q$.
  \begin{enumerate}
  \item Jika aturan terakhir adalah \olref{defn:aconvone4}, maka $P$
    berbentuk $\lambd[x][N]$ dan $Q$ berbentuk
    $\lambd[y][\Subst{N}{y}{x}]$, dengan $x \neq y$, $y \notin \FV{N}$,
    dan $\Subst{N}{y}{x}$ terdefinisi. Kita membedakan kasus menurut apakah
    $x \in \FV{N}$:
    \begin{enumerate}
    \item Jika $x \in \FV{N}$, maka:
      \begin{align*}
        \FV{\lambd[y][\Subst{N}{y}{x}]} &=
        \FV{\Subst{N}{y}{x}} \setminus \{y\} \\
        & = ((\FV{N} \setminus \{x\}) \cup \{y\}) \setminus \{y\}
         && \text{menurut \olref[sub]{thm:infv}} \\
        & = \FV{N} \setminus \{x\} \\
        & = \FV{\lambd[x][N]}
      \end{align*}
    \item Jika $x \notin \FV{N}$, maka:
      \begin{align*}
        \FV{\lambd[y][\Subst{N}{y}{x}]}
        & = \FV{\Subst{N}{y}{x}} \setminus \{y\} \\
        & = \FV{N} \setminus \{x\}
         && \text{menurut \olref[sub]{thm:notinfv}} \\
        & = \FV{\lambd[x][N]}.
      \end{align*}
    \end{enumerate}
  \item Ketiga kasus lainnya diserahkan sebagai latihan.
  \end{enumerate}
\end{proof}

\begin{prob}
  Lengkapi pembuktian \olref[lam][syn][alp]{lem:fv-one}.
\end{prob}

\begin{lem}\ollabel{lem:inv}
  Jika $P \aconvone Q$, maka $Q \aconvone P$.
\end{lem}

\begin{proof}
  Induksi pada !!{derivation} $P \aconvone Q$.
  \begin{enumerate}
  \item Jika aturan terakhir adalah \olref{defn:aconvone4}, maka $P$
    berbentuk $\lambd[x][N]$ dan $Q$ berbentuk
    $\lambd[y][\Subst{N}{y}{x}]$, dengan $x \neq y$, $y \notin \FV{N}$,
    dan $\Subst{N}{y}{x}$ terdefinisi. Pertama, kita mempunyai $y \notin
    \FV{\Subst{N}{y}{x}}$ menurut \olref[sub]{thm:clr}. Menurut
    \olref[sub]{thm:inv}, $\Subst{\Subst{N}{y}{x}}{x}{y}$ tidak hanya
    terdefinisi, tetapi juga sama dengan~$N$. Maka, berdasarkan
    \olref{defn:aconvone4}, kita mempunyai
    $\lambd[y][\Subst{N}{y}{x}] \aconvone
    \lambd[x][\Subst{\Subst{N}{y}{x}}{x}{y}] = \lambd[x][N]$.
  \end{enumerate}
\end{proof}

\begin{prob}
  Lengkapi pembuktian \olref[lam][syn][alp]{lem:inv}.
\end{prob}

\begin{thm}
  Konversi-$\alpha$ merupakan relasi ekuivalensi pada suku-suku; yakni, relasi
  itu refleksif, simetris, dan transitif.
\end{thm}

\begin{proof}
  \begin{enumerate}
  \item Untuk setiap suku $M$, $M$ dapat diubah menjadi $M$ dengan
    \emph{nol} kali pengubahan variabel terikat.
  \item Jika $P$ dapat dikonversi-$\alpha$ menjadi $Q$ melalui serangkaian
    pengubahan variabel terikat, maka dari $Q$ kita cukup membalik pengubahan
    tersebut (berdasarkan \olref{lem:inv}) dalam urutan terbalik untuk
    memperoleh~$P$.
  \item Jika $P$ dapat dikonversi-$\alpha$ menjadi $Q$ melalui serangkaian
    pengubahan variabel terikat, dan $Q$ menjadi $R$ melalui rangkaian lain,
    maka kita dapat mengubah $P$ menjadi $R$ dengan terlebih dahulu menerapkan
    rangkaian pertama, lalu rangkaian kedua.
  \end{enumerate}
\end{proof}

Mulai sekarang, kita mengatakan bahwa $M$ dan $N$ \emph{$\alpha$-ekuivalen},
$M \aeq N$, jika dan hanya jika $M$ dapat dikonversi-$\alpha$ menjadi $N$
(yang, sebagaimana baru saja kita tunjukkan, berlaku jika dan hanya jika $N$
dapat dikonversi-$\alpha$ menjadi~$M$).

\begin{thm}\ollabel{thm:fv}
  Jika $M \aeq N$, maka $\FV{M} = \FV{N}$.
\end{thm}

\begin{proof}
  Langsung dari \olref{lem:fv-one}.
\end{proof}

\begin{lem}\ollabel{lem:sub:R}
  Jika $R \aeq R'$ dan $\Subst{M}{R}{y}$ terdefinisi, maka
  $\Subst{M}{R'}{y}$ terdefinisi dan $\alpha$-ekuivalen dengan
  $\Subst{M}{R}{y}$.
\end{lem}

\begin{proof}
  Latihan.
\end{proof}

\begin{prob}
  Buktikan \olref[lam][syn][alp]{lem:sub:R}.
\end{prob}

Ingat bahwa dalam \olref[sub]{sec}, substitusi tidak terdefinisi dalam
beberapa kasus. Namun, dengan menggunakan konversi-$\alpha$ pada suku-suku,
kita dapat membuat substitusi selalu terdefinisi dengan mengganti nama
variabel terikat. Hasilnya mempertahankan ekuivalensi-$\alpha$, sebagaimana
ditunjukkan dalam teorema berikut:

\begin{thm}\ollabel{thm:sub}
  Untuk sebarang $M$, $R$, dan~$y$, terdapat $M'$ sedemikian sehingga
  $M \aeq M'$ dan $\Subst{M'}{R}{y}$ terdefinisi. Lebih lanjut, jika terdapat
  pasangan lain $M'' \aeq M$ dan $R''$ dengan $\Subst{M''}{R''}{y}$
  terdefinisi dan $R'' \aeq R$, maka
  $\Subst{M'}{R}{y} \aeq \Subst{M''}{R''}{y}$.
\end{thm}

\begin{proof}
  Dengan induksi pada pembentukan $M$:
  \begin{enumerate}
  \item $M$ adalah variabel~$z$: latihan.
  \item Andaikan $M$ berbentuk $\lambd[x][N]$. Pilih suatu variabel $z$ yang
    berbeda dari $x$ dan $y$, serta sedemikian sehingga $z \notin \FV{N}$
    dan $z \notin \FV{R}$. Berdasarkan hipotesis induksi, terdapat $N'$
    sedemikian sehingga $N' \aeq N$ dan $\Subst{N'}{z}{x}$ terdefinisi. Maka
    $\lambd[x][N] \aeq \lambd[x][N']$ juga, berdasarkan
    \olref{defn:aconvone}\olref{defn:aconvone1}. Sekarang
    $\lambd[x][N'] \aeq \lambd[z][\Subst{N'}{z}{x}]$ berdasarkan
    \olref{defn:aconvone}\olref{defn:aconvone4}. Kita dapat melakukan hal ini
    karena $z \ne x$, $z \notin \FV{N'}$, dan $\Subst{N'}{z}{x}$
    terdefinisi. Terakhir,
    $\Subst{\lambd[z][\Subst{N'}{z}{x}]}{R}{y}$ terdefinisi, karena
    $z \neq y$ dan $z \notin \FV{R}$.

    Lebih lanjut, jika terdapat $N''$ dan $R''$ lain yang memenuhi
    syarat-syarat yang sama,
    \begin{multline*}
      \Subst{(\lambd[z][\Subst{N''}{z}{x}])}{R''}{y} \\
    \begin{aligned}
      &= \lambd[z][\Subst{\Subst{N''}{z}{x}}{R''}{y}] \\
      &\aeq \lambd[z][\Subst{\Subst{N''}{z}{x}}{R}{y}] && \text{menurut
        \olref{lem:sub:R}}\\
      &\aeq \lambd[z][\Subst{\Subst{N'}{z}{x}}{R}{y}]
      && \text{menurut hipotesis induksi}\\
      &=\Subst{(\lambd[z][\Subst{N'}{z}{x}])}{R}{y}
    \end{aligned}
    \end{multline*}
    \item $M$ berbentuk $(PQ)$: latihan.
  \end{enumerate}
\end{proof}

\begin{prob}
  Lengkapi pembuktian \olref[lam][syn][alp]{thm:sub}.
\end{prob}

\begin{cor}\ollabel{cor:sub}
  Untuk sebarang $M$, $R$, dan~$y$, terdapat pasangan $M'$ dan $R'$
  sedemikian sehingga $M' \aeq M$, $R \aeq R'$, dan $\Subst{M'}{R'}{y}$
  terdefinisi. Lebih lanjut, jika terdapat pasangan lain $M'' \aeq M$ dan
  $R'' \aeq R$ dengan $\Subst{M''}{R''}{y}$ terdefinisi, maka
  $\Subst{M'}{R'}{y} \aeq \Subst{M''}{R''}{y}$.
\end{cor}

\begin{proof}
  Langsung dari \olref{thm:sub}.
\end{proof}

\end{document}
